\chapter{2003:Peter Agre and Roderick MacKinnon}

\label{chap:chemistry2003}

The Nobel Prize in Chemistry for the year 2003 was awarded jointly to Peter Agre and Roderick MacKinnon.

\textbf{Motivation:}

for the discovery of water channels (one half, Peter Agre) and for structural and mechanistic studies of ion channels (the other half, Roderick MacKinnon)

\section{Peter Agre}

\subsection*{Early Life and Education}

Peter Agre was born in 1949 in Northfield, Minnesota, USA.

\subsection*{Career and Major Research}

Agre studied at Augsburg College (B.A., 1970) and earned his M.D. from the Johns Hopkins University in 1974. He was a professor at the Johns Hopkins University and later at the Johns Hopkins Bloomberg School of Public Health.

\subsection*{Nobel Prize-winning Work}

The work for which Peter Agre was awarded the Nobel Prize in Chemistry in 2003 was described by the Nobel Committee as: "for the discovery of water channels".

He discovered aquaporins, integral membrane proteins that form water channels and allow rapid water transport across cell membranes.

\subsection*{Significance and Impact}

The discovery explained how water passes through biological membranes much faster than simple diffusion would permit; aquaporins are essential in the kidneys and many tissues.

\subsection*{Later Career and Honors}

Agre continued his research at the Johns Hopkins University, where he directed the Malaria Research Institute. He shared the 2003 Nobel Prize in Chemistry.

\subsection*{Legacy}

Aquaporins are now known to be present in virtually all living organisms.

\subsection*{Modern Developments}

Agre's discovery of aquaporins, the water channels that let water cross cell membranes rapidly, solved the long-standing puzzle of water transport and founded the field of aquaporin biology. Aquaporins are now targets for diuretic and anti-cancer drugs and inspire the design of biomimetic membranes for water purification.

\subsection*{Selected Publications}

"Appearance of Water Channels in Xenopus Oocytes Expressing Red Cell CHIP28 Protein" (Science, 1992)

\clearpage

\section{Roderick MacKinnon}

\subsection*{Early Life and Education}

Roderick MacKinnon was born in 1956 in Burlington, Massachusetts, USA.

\subsection*{Career and Major Research}

MacKinnon studied at Brandeis University (B.A., 1978) and earned his M.D. from Tufts University in 1982. He joined Rockefeller University in 1996.

\subsection*{Nobel Prize-winning Work}

The work for which Roderick MacKinnon was awarded the Nobel Prize in Chemistry in 2003 was recognized for structural and mechanistic studies of ion channels.

He determined the three-dimensional structure of a potassium channel by X-ray crystallography, revealing how the selectivity filter distinguishes potassium ions from sodium ions.

\subsection*{Significance and Impact}

The structure explained the molecular basis of ion selectivity and conduction that underlies nerve signaling.

\subsection*{Later Career and Honors}

MacKinnon remained at Rockefeller University, where he is an investigator of the Howard Hughes Medical Institute. He shared the 2003 Nobel Prize in Chemistry.

\subsection*{Legacy}

His structures of potassium channels clarified how the nervous system transmits electrical signals.

\subsection*{Modern Developments}

MacKinnon's determination of the three-dimensional structures of potassium channels revealed how ion channels select and conduct ions with astonishing specificity and speed. His work explained the molecular basis of nerve and muscle signaling and underpins the pharmacology of ion-channel drugs used to treat heart disease, epilepsy, and pain.

\subsection*{Selected Publications}

"The Structure of the Potassium Channel: Molecular Basis of K+ Conduction and Selectivity" (Science, 1998)

\clearpage

\section*{Chapter Summary}

The 2003 Nobel Prize in Chemistry recognized the channels that control water and ion movement across cell membranes. Agre discovered the aquaporin water channels, and MacKinnon solved the structure of potassium channels.
